\documentclass[11pt,aspectratio=169]{beamer}
\usetheme{SimplePlus}

\usepackage[T1]{fontenc}
\usepackage[utf8]{inputenc}
\usepackage{amsmath,amssymb,mathtools}
\usepackage{booktabs}
\usepackage{array}
\usepackage{appendixnumberbeamer}
\usepackage{hyperref}

\newcommand{\E}{\mathbb{E}}
\newcommand{\Prob}{\mathbb{P}}
\newcommand{\Var}{\mathrm{Var}}
\newcommand{\Cov}{\mathrm{Cov}}
\newcommand{\ATE}{\mathrm{ATE}}
\newcommand{\ATT}{\mathrm{ATT}}
\newcommand{\argmin}{\mathop{\mathrm{arg\,min}}}
\newcommand{\indep}{\perp\!\!\!\perp}
\newcommand{\derivbutton}[1]{\vspace{0.4em}\hfill\hyperlink{app:#1}{\beamerbutton{Appendix details}}}
\newcommand{\backbutton}[1]{\vspace{0.4em}\hfill\hyperlink{main:#1}{\beamerbutton{Back to main slide}}}

\title{Applied Econometrics: Session 4}
\subtitle{Regression, Dummy Variables, Grouped Data, and Controls}
\author{PhD Intensive Course}
\institute{Lecture 4}
\date{\today}

\begin{document}

\begin{frame}
  \titlepage
\end{frame}

\begin{frame}{Why This Lecture Matters}
  Linear regression is the workhorse of applied causal inference.
  \vspace{0.4em}

  But regression has two very different meanings:
  \begin{itemize}
    \item mechanically, it is a best linear summary of conditional means;
    \item causally, it is valid only when the comparison left inside the coefficient is credible.
  \end{itemize}
  \begin{block}{Core question}
    Which variation in treatment remains after we control for other variables, and is that variation as-good-as-random?
  \end{block}
\end{frame}

\begin{frame}{Roadmap for Today}
  \tableofcontents
\end{frame}

\section{CEF, OLS, and Causal Regression}

\begin{frame}{Where We Are in the Course}
  Earlier lectures gave us three ingredients:
  \begin{itemize}
    \item potential outcomes define the causal target;
    \item random assignment removes selection bias by design;
    \item regression summarizes conditional averages and partial correlations.
  \end{itemize}
  \vspace{0.3em}
  Today we put them together:
  \begin{block}{Lecture 4 message}
    Regression adjustment is powerful when it compares treated and untreated units that are comparable after conditioning. It is dangerous when controls create artificial comparisons.
  \end{block}
\end{frame}

\begin{frame}{Potential Outcomes Reminder}
  \hypertarget{main:ate-reminder}{}
  For each unit \(i\):
  \[
    Y_i = Y_{0i} + T_i(Y_{1i}-Y_{0i}), \qquad T_i\in\{0,1\}.
  \]
  The average treatment effect is
  \[
    \ATE = \E[Y_{1i}-Y_{0i}].
  \]
  A simple observed difference satisfies
  \[
    \E[Y_i\mid T_i=1]-\E[Y_i\mid T_i=0]
    =
    \ATT
    +
    \left\{\E[Y_{0i}\mid T_i=1]-\E[Y_{0i}\mid T_i=0]\right\}.
  \]
  \begin{itemize}
    \item The second term is selection bias.
    \item Regression is useful only if it helps remove or avoid this term.
  \end{itemize}
  \derivbutton{ate-reminder}
\end{frame}

\begin{frame}{Regression with a Treatment Dummy}
  \hypertarget{main:dummy-diff}{}
  In a randomized experiment, estimate
  \[
    Y_i=\alpha+\tau T_i+u_i.
  \]
  Since \(T_i\) is a dummy:
  \[
    \E[Y_i\mid T_i=0]=\alpha,\qquad
    \E[Y_i\mid T_i=1]=\alpha+\tau.
  \]
  Therefore
  \[
    \tau=\E[Y_i\mid T_i=1]-\E[Y_i\mid T_i=0].
  \]
  \begin{block}{Interpretation}
    With a binary treatment and an intercept, OLS is just a difference in means written in regression language.
  \end{block}
  \derivbutton{dummy-diff}
\end{frame}

\begin{frame}{The Conditional Expectation Function}
  \hypertarget{main:cef}{}
  For outcome \(Y_i\) and covariates \(X_i\), define
  \[
    m(x)=\E[Y_i\mid X_i=x].
  \]
  \begin{itemize}
    \item The CEF is the average outcome among units with the same covariates.
    \item It is the best unrestricted prediction of \(Y_i\) using \(X_i\).
    \item It can be nonlinear, discontinuous, or defined by many cells.
    \item It is not automatically causal.
  \end{itemize}
  \begin{block}{Beginner intuition}
    If \(X\) is ``years of education'', the CEF asks: among people with exactly \(x\) years, what is the average wage?
  \end{block}
  \derivbutton{cef}
\end{frame}

\begin{frame}{Population OLS}
  \hypertarget{main:ols-normal}{}
  Regression chooses the best linear approximation:
  \[
    \beta
    =
    \argmin_b \E\left[(Y_i-X_i'b)^2\right].
  \]
  The first-order condition gives
  \[
    \E[X_i(Y_i-X_i'\beta)]=0.
  \]
  If \(\E[X_iX_i']\) is invertible,
  \[
    \beta=\E[X_iX_i']^{-1}\E[X_iY_i].
  \]
  \begin{itemize}
    \item Orthogonality is a property of the projection.
    \item It does not say the treatment is randomly assigned.
  \end{itemize}
  \derivbutton{ols-normal}
\end{frame}

\begin{frame}{Regression Anatomy: What Multiple Regression Does}
  \hypertarget{main:fwl}{}
  Suppose the regression is
  \[
    Y_i=\alpha+\tau T_i+X_i'\gamma+u_i.
  \]
  Let \(\tilde T_i\) be the residual from regressing \(T_i\) on \(X_i\) and a constant.
  Then the coefficient on treatment is
  \[
    \tau
    =
    \frac{\Cov(Y_i,\tilde T_i)}{\Var(\tilde T_i)}.
  \]
  \begin{block}{Intuition}
    Multiple regression first removes from treatment the part predictable from controls. Then it asks whether the remaining treatment variation predicts the outcome.
  \end{block}
  \derivbutton{fwl}
\end{frame}

\begin{frame}{When Regression Adjustment Is Causal}
  \hypertarget{main:cia-reg}{}
  Regression adjustment needs a causal assumption, not only algebra:
  \[
    (Y_{0i},Y_{1i})\indep T_i\mid X_i.
  \]
  This is conditional independence, or selection on observables.
  \vspace{0.3em}

  Under this assumption,
  \[
    \E[Y_i\mid T_i=1,X_i=x]-\E[Y_i\mid T_i=0,X_i=x]
    =
    \E[Y_{1i}-Y_{0i}\mid X_i=x].
  \]
  \begin{block}{Meaning}
    Within each \(X\)-cell, treated and untreated units are comparable in potential outcomes.
  \end{block}
  \derivbutton{cia-reg}
\end{frame}

\begin{frame}{Regression Adjustment in Words}
  A regression with controls tries to imitate a blocked experiment:
  \begin{enumerate}
    \item compare treated and untreated units with the same observed covariates;
    \item estimate many within-cell treatment differences;
    \item average those differences using weights implied by the regression.
  \end{enumerate}
  \vspace{0.3em}
  \begin{block}{What regression cannot do}
    It cannot fix unobserved confounders, bad measurement, lack of overlap, or post-treatment conditioning mistakes.
  \end{block}
\end{frame}

\section{Omitted-Variable Bias}

\begin{frame}{Why Omitting Confounders Biases Regression}
  A confounder is a variable that affects both treatment and outcome.
  \[
    A_i \rightarrow T_i,\qquad A_i \rightarrow Y_i.
  \]
  Example: ability in a wage-schooling regression.
  \begin{itemize}
    \item Ability raises wages directly.
    \item Ability also predicts years of education.
    \item A short regression of wage on education mixes education and ability.
  \end{itemize}
  \begin{block}{Plain language}
    If high-education people would have earned more even without extra education, the education coefficient is too large.
  \end{block}
\end{frame}

\begin{frame}{The Omitted-Variable Bias Formula}
  \hypertarget{main:ovb}{}
  Suppose the long regression is
  \[
    Y_i=\alpha+\tau T_i+\beta A_i+u_i,
  \]
  but we estimate the short regression \(Y_i\) on \(T_i\) only.
  Then
  \[
    \tau_{\text{short}}
    =
    \tau+\beta\delta,
  \]
  where
  \[
    \delta = \frac{\Cov(A_i,T_i)}{\Var(T_i)}
  \]
  is the slope from regressing the omitted variable \(A_i\) on treatment \(T_i\).
  \begin{block}{Memorize}
    Short coefficient \(=\) long coefficient \(+\) effect of omitted variable on outcome \(\times\) relationship between omitted and included variable.
  \end{block}
  \derivbutton{ovb}
\end{frame}

\begin{frame}{Sign Intuition for OVB}
  \begin{center}
    \begin{tabular}{@{}ccc@{}}
      \toprule
      Effect of omitted \(A\) on \(Y\) & Corr. of \(A\) with \(T\) & Bias direction \\
      \midrule
      Positive & Positive & Upward \\
      Positive & Negative & Downward \\
      Negative & Positive & Downward \\
      Negative & Negative & Upward \\
      \bottomrule
    \end{tabular}
  \end{center}
  \vspace{0.4em}
  \begin{block}{Exam habit}
    State both signs. ``The coefficient is biased'' is incomplete; the sign depends on both arrows.
  \end{block}
\end{frame}

\begin{frame}{OVB Simulation: Education and Wages}
  The script simulates latent ability and family background.
  \vspace{0.3em}

  \begin{center}
    \begin{tabular}{@{}lcc@{}}
      \toprule
      Specification & Education coefficient & Interpretation \\
      \midrule
      Short: no controls & 0.204 & badly upward biased \\
      Add noisy proxies & 0.143 & bias falls but remains \\
      Oracle: true confounders & 0.078 & close to true 0.080 \\
      \bottomrule
    \end{tabular}
  \end{center}
  \vspace{0.3em}
  \begin{itemize}
    \item The short regression attributes ability and family background to education.
    \item Imperfect proxies help, but they do not fully reconstruct the experiment.
    \item The oracle regression is not available in real work; it is a teaching benchmark.
  \end{itemize}
\end{frame}

\begin{frame}{FWL Simulation: What Changed After Controls?}
  In the same simulation, the education coefficient with proxies is \(0.143\).
  \vspace{0.4em}

  The Frisch--Waugh--Lovell calculation:
  \begin{enumerate}
    \item regress education on test score and parent index;
    \item save residualized education \(\tilde{Educ}_i\);
    \item regress log wage on \(\tilde{Educ}_i\).
  \end{enumerate}
  gives \(0.143\) again.
  \vspace{0.4em}
  \begin{block}{Interpretation}
    The controlled coefficient uses only the part of education not linearly predicted by the included controls.
  \end{block}
\end{frame}

\begin{frame}{Controls Are Not a Ritual}
  Do not add controls because the dataset contains them.
  \vspace{0.3em}

  For each proposed control, ask:
  \begin{enumerate}
    \item Is it determined before treatment?
    \item Does it help block a back-door path from treatment to outcome?
    \item Does it mainly improve outcome prediction without changing treatment assignment?
    \item Could treatment itself affect it?
    \item Could it be a common effect of treatment and an outcome determinant?
  \end{enumerate}
  \begin{block}{Practical rule}
    A control list is part of the research design. It needs an argument, not only a software formula.
  \end{block}
\end{frame}

\section{Dummy Variables and Grouped Regression}

\begin{frame}{Binary Dummy Variables}
  A dummy variable is a \(0/1\) indicator.
  \[
    Female_i =
    \begin{cases}
      1 & \text{if unit } i \text{ is female},\\
      0 & \text{otherwise}.
    \end{cases}
  \]
  In
  \[
    Y_i=\alpha+\beta Female_i+u_i,
  \]
  we have
  \[
    \alpha=\E[Y_i\mid Female_i=0],
    \qquad
    \beta=\E[Y_i\mid Female_i=1]-\E[Y_i\mid Female_i=0].
  \]
  \begin{block}{Important}
    A dummy coefficient is a difference relative to a base group.
  \end{block}
\end{frame}

\begin{frame}{The Dummy-Variable Trap}
  Suppose education has categories \(8,9,\ldots,18\).
  \vspace{0.3em}

  If we include an intercept and a dummy for every category, the columns add up to the intercept:
  \[
    1 = D_{8i}+D_{9i}+\cdots+D_{18i}.
  \]
  Then \(X'X\) is singular.
  \vspace{0.3em}

  Standard solution:
  \begin{itemize}
    \item keep the intercept;
    \item omit one category as the base group;
    \item interpret every dummy relative to that omitted category.
  \end{itemize}
\end{frame}

\begin{frame}{Many Dummies: A Saturated Model}
  \hypertarget{main:dummy-means}{}
  With categories \(g=1,\ldots,G\), estimate
  \[
    Y_i=\alpha+\beta_2D_{2i}+\cdots+\beta_GD_{Gi}+u_i,
  \]
  where group 1 is the base group.
  \vspace{0.3em}

  OLS fitted values are group means:
  \[
    \widehat{Y}_i=
    \begin{cases}
      \bar Y_1 & \text{if } i\in g=1,\\
      \bar Y_g & \text{if } i\in g>1.
    \end{cases}
  \]
  \begin{block}{Interpretation}
    A saturated dummy regression is a table of conditional means written as a regression.
  \end{block}
  \derivbutton{dummy-means}
\end{frame}

\begin{frame}{Dummy Simulation: Flexible Education Effects}
  Simulated wages have a nonlinear education schedule.
  \vspace{0.3em}

  \begin{center}
    \begin{tabular}{@{}lcc@{}}
      \toprule
      Model & What it imposes & Example output \\
      \midrule
      Linear \(Y\) on education & one constant slope & slope \(=1.691\) \\
      Education dummies & one mean per year & fitted means equal cell means \\
      \bottomrule
    \end{tabular}
  \end{center}
  \vspace{0.3em}
  \begin{block}{Why use dummies?}
    Dummies avoid forcing a straight line when the CEF is naturally step-shaped or nonlinear.
  \end{block}
\end{frame}

\begin{frame}{Treatment Dummies and Interactions}
  Suppose \(T_i\) is treatment and \(G_i\) is a group dummy.
  \[
    Y_i=\alpha+\tau T_i+\lambda G_i+u_i.
  \]
  This allows different baseline means by group, but the treatment effect is forced to be the same.
  \vspace{0.3em}

  To allow different treatment effects:
  \[
    Y_i=\alpha+\tau T_i+\lambda G_i+\theta(T_iG_i)+u_i.
  \]
  Then
  \[
    \text{effect for }G=0:\tau,\qquad
    \text{effect for }G=1:\tau+\theta.
  \]
  \begin{block}{Student warning}
    ``Controlling for group'' is not the same as allowing the effect to differ by group.
  \end{block}
\end{frame}

\begin{frame}{Grouped Data: Why Weights Matter}
  \hypertarget{main:grouped}{}
  Sometimes data are aggregated:
  \[
    (\bar Y_g,\bar X_g,n_g),\qquad g=1,\ldots,G.
  \]
  If \(X_i\) is constant within group \(g\), individual OLS on
  \[
    Y_i=\alpha+\beta X_i+u_i
  \]
  can be recovered from grouped means by weighted least squares:
  \[
    \min_{\alpha,\beta}\sum_g n_g(\bar Y_g-\alpha-\beta \bar X_g)^2.
  \]
  \begin{block}{Reason}
    A group with 800 observations should count more than a group with 8 observations.
  \end{block}
  \derivbutton{grouped}
\end{frame}

\begin{frame}{Grouped Regression Simulation}
  In the simulated education-wage data:
  \vspace{0.3em}

  \begin{center}
    \begin{tabular}{@{}lcc@{}}
      \toprule
      Estimator & Intercept & Education slope \\
      \midrule
      Individual OLS & -2.247 & 1.691 \\
      Grouped WLS with counts & -2.247 & 1.691 \\
      Grouped unweighted OLS & -4.224 & 1.888 \\
      \bottomrule
    \end{tabular}
  \end{center}
  \vspace{0.3em}
  \begin{itemize}
    \item Weighted grouped regression reproduces the individual-data coefficient here.
    \item Unweighted grouped regression treats rare and common education levels equally.
  \end{itemize}
\end{frame}

\begin{frame}{When Grouped Regression Is Safe}
  Grouped regression is safe when:
  \begin{itemize}
    \item the group means correspond to the same units and population;
    \item the regressors used in the model are available as group means;
    \item weights are the number of observations behind each mean;
    \item aggregation does not hide within-group confounding needed for identification.
  \end{itemize}
  \begin{block}{Common mistake}
    Regressing totals on totals often changes the question. For grouped OLS, use means with counts as weights.
  \end{block}
\end{frame}

\begin{frame}{What Dummies and Groups Teach Us}
  Dummies and grouped data make a larger point:
  \begin{itemize}
    \item regression can reproduce cell means exactly when the model is saturated;
    \item parametric regressions impose structure across cells;
    \item weights decide which cells matter most;
    \item causal interpretation still depends on comparability within cells.
  \end{itemize}
  \begin{block}{Bridge to controls}
    More controls create more cells. This can help identification, but it can also create sparse cells, bad overlap, and bad-control mistakes.
  \end{block}
\end{frame}

\section{Beyond Confounders}

\begin{frame}{Not Every Control Has the Same Role}
  Controls can play different roles:
  \begin{center}
    \begin{tabular}{@{}lll@{}}
      \toprule
      Type & Graphical role & Usual advice \\
      \midrule
      Confounder & \(X\rightarrow T,\ X\rightarrow Y\) & include \\
      Outcome predictor & \(X\rightarrow Y\) only & often include for precision \\
      Treatment predictor & \(X\rightarrow T\) only & usually avoid if not needed \\
      Mediator & \(T\rightarrow M\rightarrow Y\) & exclude for total effect \\
      Collider & \(T\rightarrow C\leftarrow U\rightarrow Y\) & exclude \\
      \bottomrule
    \end{tabular}
  \end{center}
  \vspace{0.3em}
  The same word, ``control'', hides very different causal objects.
\end{frame}

\begin{frame}{Outcome Predictors Can Improve Precision}
  Suppose treatment is randomized:
  \[
    Y_i=\alpha+\tau T_i+\gamma X_i+\varepsilon_i,
  \]
  where \(X_i\) predicts \(Y_i\) but not \(T_i\).
  \vspace{0.3em}

  Including \(X_i\):
  \begin{itemize}
    \item does not fix bias, because there was no bias from random assignment;
    \item can reduce residual variance;
    \item often makes the standard error for \(\hat\tau\) smaller.
  \end{itemize}
  \begin{block}{Example}
    Baseline test scores in an education experiment often improve precision for treatment effects on final test scores.
  \end{block}
\end{frame}

\begin{frame}{Treatment Predictors Can Hurt Precision}
  A variable can predict treatment but add little outcome information after other controls.
  \vspace{0.3em}

  By FWL,
  \[
    \hat\tau=\frac{\Cov(Y_i,\tilde T_i)}{\Var(\tilde T_i)}.
  \]
  If a new control predicts \(T_i\), then \(\Var(\tilde T_i)\) can shrink.
  \vspace{0.3em}

  \begin{block}{Consequence}
    Less residual treatment variation means a noisier estimate, unless the control also helps explain the outcome enough to compensate.
  \end{block}
  \derivbutton{precision}
\end{frame}

\begin{frame}{Bad Controls: Mediators}
  \hypertarget{main:mediator}{}
  A mediator lies on the causal path:
  \[
    T_i \rightarrow M_i \rightarrow Y_i.
  \]
  If the target is the total effect of \(T_i\) on \(Y_i\), do not control for \(M_i\).
  \vspace{0.3em}

  Example: email treatment \(\rightarrow\) customer opens email \(\rightarrow\) customer pays.
  \begin{itemize}
    \item Total effect asks: what is the effect of sending the email?
    \item Controlling for opening asks: what is the effect of sending the email holding opening fixed?
    \item The second question removes part of the mechanism.
  \end{itemize}
  \derivbutton{mediator}
\end{frame}

\begin{frame}{Mediator Simulation}
  Treatment is randomized. It increases agreement, and agreement increases payment.
  \vspace{0.3em}

  \begin{center}
    \small
    \begin{tabular}{@{}lcc@{}}
      \toprule
      Regression & Treatment coefficient & Meaning \\
      \midrule
      Payment on treatment & 8.636 & total effect estimate \\
      Payment on treatment + agreement & 2.211 & direct effect \\
      \bottomrule
    \end{tabular}
  \end{center}
  \vspace{0.3em}
  \begin{block}{Lesson}
    A control can make a randomized experiment answer the wrong question.
  \end{block}
\end{frame}

\begin{frame}{Bad Controls: Colliders and Selection}
  \hypertarget{main:collider}{}
  A collider is a common effect:
  \[
    T_i \rightarrow C_i \leftarrow U_i \rightarrow Y_i.
  \]
  If we condition on \(C_i\), treatment and the unobserved outcome determinant \(U_i\) become associated.
  \vspace{0.3em}

  \begin{block}{Intuition}
    Among units with the same \(C\), a high value of treatment must be offset by a lower value of \(U\). This creates a non-causal comparison.
  \end{block}
  \derivbutton{collider}
\end{frame}

\begin{frame}{Collider Simulation}
  Treatment is randomized and the true effect is 6.
  \vspace{0.3em}

  \begin{center}
    \begin{tabular}{@{}lcc@{}}
      \toprule
      Regression & Treatment coefficient & Interpretation \\
      \midrule
      Outcome on treatment & 5.999 & randomization works \\
      Outcome on treatment + collider & 1.700 & conditioning creates bias \\
      \bottomrule
    \end{tabular}
  \end{center}
  \vspace{0.3em}
  \begin{itemize}
    \item The collider is partly caused by treatment.
    \item It is also partly caused by an unobserved determinant of the outcome.
    \item Holding it fixed opens a non-causal path from treatment to outcome.
  \end{itemize}
\end{frame}

\begin{frame}{Selection on Outcomes Is Also a Bad Control Problem}
  Many applied mistakes look like:
  \[
    \text{estimate treatment effect only among units with }Y_i>0.
  \]
  \vspace{0.3em}
  If treatment affects whether \(Y_i>0\), this selects on a post-treatment outcome.
  \begin{itemize}
    \item Example: marketing campaign affects whether customers buy anything.
    \item Dropping zero spenders keeps different latent customer types in treatment and control.
    \item Random assignment before selection does not rescue the selected comparison.
  \end{itemize}
  \begin{block}{Safer target}
    Estimate the effect on total spending for everyone, then separately describe participation if that is policy-relevant.
  \end{block}
\end{frame}

\begin{frame}{Practical Control-Selection Rules}
  Before adding \(X\), classify it:
  \begin{enumerate}
    \item \textbf{Pre-treatment confounder:} include if it blocks selection bias.
    \item \textbf{Pre-treatment outcome predictor:} include if it improves precision and does not damage overlap.
    \item \textbf{Treatment predictor only:} be cautious; it may reduce useful treatment variation.
    \item \textbf{Mediator:} exclude for total effects; include only for a direct-effect question with extra assumptions.
    \item \textbf{Collider or selection variable:} exclude.
  \end{enumerate}
  \begin{block}{Workflow}
    Draw the timing first. If a variable is measured after treatment, treat it as suspicious until proven otherwise.
  \end{block}
\end{frame}

\begin{frame}{A Good Regression Table Explains the Design}
  A regression table should not be a list of increasingly large kitchen-sink specifications.
  \vspace{0.3em}

  Better sequence:
  \begin{enumerate}
    \item show raw treatment-control difference;
    \item add pre-treatment confounders needed for identification;
    \item add precision controls separately if useful;
    \item avoid post-treatment variables in the main total-effect specification;
    \item explain why the preferred column matches the causal question.
  \end{enumerate}
  \begin{block}{Student habit}
    Label columns by identification logic, not by ``more controls''.
  \end{block}
\end{frame}

\section{Synthesis}

\begin{frame}{What Regression Is Doing}
  \begin{itemize}
    \item A treatment dummy without controls is a difference in means.
    \item OLS approximates the CEF with a linear function.
    \item Multiple regression uses residual treatment variation after partialling out controls.
    \item Dummy regressions can reproduce group means exactly.
    \item Grouped regression needs weights when group sizes differ.
  \end{itemize}
\end{frame}

\begin{frame}{When Regression Is Credible}
  Regression is credible when:
  \begin{itemize}
    \item the target estimand is clear;
    \item controls are pre-treatment and justified by the assignment process;
    \item there is enough overlap after conditioning;
    \item the functional form is not doing hidden extrapolation;
    \item post-treatment variables are kept out of total-effect regressions.
  \end{itemize}
  \begin{block}{One-sentence summary}
    Regression is not the research design; it is the estimator that implements the design.
  \end{block}
\end{frame}

\begin{frame}{Python Examples for This Lecture}
  The self-contained script is:
  \[
    \texttt{lectures/code/lecture-4-regression-controls.py}
  \]
  It simulates:
  \begin{itemize}
    \item omitted-variable bias and the FWL theorem;
    \item dummy variables and grouped weighted regression;
    \item bad controls through mediators and colliders.
  \end{itemize}
  \begin{block}{Use in class}
    Ask students to change the signs and strengths of the simulated relationships, then predict how the coefficients should move before running the code.
  \end{block}
\end{frame}

\begin{frame}{Sources and References}
  Main source chapters:
  \begin{itemize}
    \item Matheus Facure Alves, \href{https://matheusfacure.github.io/python-causality-handbook/05-The-Unreasonable-Effectiveness-of-Linear-Regression.html}{Chapter 05: The Unreasonable Effectiveness of Linear Regression}.
    \item Matheus Facure Alves, \href{https://matheusfacure.github.io/python-causality-handbook/06-Grouped-and-Dummy-Regression.html}{Chapter 06: Grouped and Dummy Regression}.
    \item Matheus Facure Alves, \href{https://matheusfacure.github.io/python-causality-handbook/07-Beyond-Confounders.html}{Chapter 07: Beyond Confounders}.
  \end{itemize}
  Background econometrics:
  \begin{itemize}
    \item Angrist and Pischke, \emph{Mostly Harmless Econometrics}.
    \item Angrist and Pischke, \emph{Mastering 'Metrics}.
    \item Hernan and Robins, \emph{Causal Inference: What If}.
  \end{itemize}
\end{frame}

\begin{frame}{What Students Should Be Able to Do}
  After this lecture, students should be able to:
  \begin{itemize}
    \item explain why a treatment dummy coefficient equals a mean difference;
    \item use FWL intuition to interpret controlled regressions;
    \item derive and sign omitted-variable bias;
    \item interpret dummy coefficients relative to a base group;
    \item explain why grouped regressions need weights;
    \item distinguish confounders, precision controls, mediators, and colliders.
  \end{itemize}
\end{frame}

\appendix

\begin{frame}{Appendix Map}
  \begin{itemize}
    \item A1: Selection-bias decomposition reminder.
    \item A2: Dummy-regression difference in means.
    \item A3: CEF foundations.
    \item A4: OLS normal equations.
    \item A5: Frisch--Waugh--Lovell theorem.
    \item A6: Conditional independence and regression adjustment.
    \item A7: Omitted-variable bias formula.
    \item A8: Dummy regressions and group means.
    \item A9: Grouped weighted OLS.
    \item A10: Controls and precision.
    \item A11: Mediators and total effects.
    \item A12: Colliders and induced association.
  \end{itemize}
\end{frame}

\begin{frame}{Appendix A1: Selection-Bias Decomposition}
  \hypertarget{app:ate-reminder}{}
  Start with
  \[
    \E[Y_i\mid T_i=1]-\E[Y_i\mid T_i=0]
    =
    \E[Y_{1i}\mid T_i=1]-\E[Y_{0i}\mid T_i=0].
  \]
  Add and subtract \(\E[Y_{0i}\mid T_i=1]\):
  \[
    \E[Y_{1i}-Y_{0i}\mid T_i=1]
    +
    \left\{\E[Y_{0i}\mid T_i=1]-\E[Y_{0i}\mid T_i=0]\right\}.
  \]
  The first term is \(\ATT\). The second term is selection bias.
  \backbutton{ate-reminder}
\end{frame}

\begin{frame}{Appendix A2: Treatment Dummy Equals Difference in Means}
  \hypertarget{app:dummy-diff}{}
  Regression:
  \[
    Y_i=\alpha+\tau T_i+u_i,\qquad T_i\in\{0,1\}.
  \]
  The normal equations imply residuals average to zero within each treatment group:
  \[
    \sum_{T_i=0}(Y_i-\hat\alpha)=0,
    \qquad
    \sum_{T_i=1}(Y_i-\hat\alpha-\hat\tau)=0.
  \]
  Therefore
  \[
    \hat\alpha=\bar Y_0,
    \qquad
    \hat\alpha+\hat\tau=\bar Y_1,
  \]
  so
  \[
    \hat\tau=\bar Y_1-\bar Y_0.
  \]
  \backbutton{dummy-diff}
\end{frame}

\begin{frame}{Appendix A3: CEF Foundations}
  \hypertarget{app:cef}{}
  Define
  \[
    m(X_i)=\E[Y_i\mid X_i],\qquad \varepsilon_i=Y_i-m(X_i).
  \]
  Then
  \[
    \E[\varepsilon_i\mid X_i]
    =
    \E[Y_i-m(X_i)\mid X_i]
    =
    \E[Y_i\mid X_i]-m(X_i)
    =
    0.
  \]
  Hence \(\varepsilon_i\) is uncorrelated with any function \(g(X_i)\):
  \[
    \E[g(X_i)\varepsilon_i]
    =
    \E\{g(X_i)\E[\varepsilon_i\mid X_i]\}
    =
    0.
  \]
  \backbutton{cef}
\end{frame}

\begin{frame}{Appendix A4: OLS Normal Equations}
  \hypertarget{app:ols-normal}{}
  Minimize
  \[
    Q(b)=\E[(Y_i-X_i'b)^2].
  \]
  Differentiate:
  \[
    \frac{\partial Q(b)}{\partial b}
    =
    -2\E[X_i(Y_i-X_i'b)].
  \]
  At the minimum:
  \[
    \E[X_i(Y_i-X_i'\beta)]=0.
  \]
  Rearranging:
  \[
    \E[X_iY_i]=\E[X_iX_i']\beta,
    \qquad
    \beta=\E[X_iX_i']^{-1}\E[X_iY_i].
  \]
  \backbutton{ols-normal}
\end{frame}

\begin{frame}{Appendix A5: Frisch--Waugh--Lovell}
  \hypertarget{app:fwl}{}
  Let \(M_X\) be the residual-maker after projecting on controls \(X\):
  \[
    M_X=I-X(X'X)^{-1}X'.
  \]
  In
  \[
    Y=T\tau+X\gamma+u,
  \]
  premultiply by \(M_X\):
  \[
    M_XY=M_XT\tau+M_XX\gamma+M_Xu.
  \]
  Since \(M_XX=0\),
  \[
    \tilde Y=\tilde T\tau+\tilde u.
  \]
  Thus
  \[
    \hat\tau=(\tilde T'\tilde T)^{-1}\tilde T'\tilde Y.
  \]
  With a single treatment variable, this is \(\Cov(Y,\tilde T)/\Var(\tilde T)\).
  \backbutton{fwl}
\end{frame}

\begin{frame}{Appendix A6: Conditional Independence}
  \hypertarget{app:cia-reg}{}
  If
  \[
    (Y_{0i},Y_{1i})\indep T_i\mid X_i,
  \]
  then
  \[
    \E[Y_i\mid T_i=1,X_i=x]
    =
    \E[Y_{1i}\mid T_i=1,X_i=x]
    =
    \E[Y_{1i}\mid X_i=x],
  \]
  and
  \[
    \E[Y_i\mid T_i=0,X_i=x]
    =
    \E[Y_{0i}\mid X_i=x].
  \]
  Subtracting gives
  \[
    \E[Y_i\mid T_i=1,X_i=x]-\E[Y_i\mid T_i=0,X_i=x]
    =
    \E[Y_{1i}-Y_{0i}\mid X_i=x].
  \]
  \backbutton{cia-reg}
\end{frame}

\begin{frame}{Appendix A7: Omitted-Variable Bias Formula}
  \hypertarget{app:ovb}{}
  True long model:
  \[
    Y_i=\alpha+\tau T_i+\beta A_i+u_i,
    \qquad \Cov(T_i,u_i)=0.
  \]
  Short-regression slope:
  \[
    \tau_{\text{short}}=\frac{\Cov(Y_i,T_i)}{\Var(T_i)}.
  \]
  Substitute the true model:
  \[
    \tau_{\text{short}}
    =
    \frac{\Cov(\alpha+\tau T_i+\beta A_i+u_i,T_i)}{\Var(T_i)}.
  \]
  Using covariance rules:
  \[
    \tau_{\text{short}}
    =
    \tau
    +
    \beta\frac{\Cov(A_i,T_i)}{\Var(T_i)}
    =
    \tau+\beta\delta.
  \]
  \backbutton{ovb}
\end{frame}

\begin{frame}{Appendix A8: Dummy Regression Gives Group Means}
  \hypertarget{app:dummy-means}{}
  \small
  For groups \(g=1,\ldots,G\), estimate a saturated model with one omitted base group:
  \[
    Y_i=\alpha+\sum_{g=2}^G\beta_gD_{gi}+u_i.
  \]
  Fitted values are constant within group. OLS chooses constants \(c_g\) to minimize
  \[
    \sum_{g=1}^G\sum_{i\in g}(Y_i-c_g)^2.
  \]
  Each \(c_g\) solves
  \[
    \frac{\partial}{\partial c_g}\sum_{i\in g}(Y_i-c_g)^2
    =
    -2\sum_{i\in g}(Y_i-c_g)=0,
  \]
  so
  \[
    c_g=\bar Y_g.
  \]
  \backbutton{dummy-means}
\end{frame}

\begin{frame}{Appendix A9: Grouped Weighted OLS}
  \hypertarget{app:grouped}{}
  Individual least squares with groups:
  \[
    \sum_i(Y_i-\alpha-\beta X_i)^2
    =
    \sum_g\sum_{i\in g}(Y_i-\alpha-\beta X_g)^2.
  \]
  Decompose \(Y_i=\bar Y_g+(Y_i-\bar Y_g)\):
  \[
    \sum_{i\in g}(Y_i-\alpha-\beta X_g)^2
    =
    n_g(\bar Y_g-\alpha-\beta X_g)^2
    +
    \sum_{i\in g}(Y_i-\bar Y_g)^2.
  \]
  The second term does not depend on \(\alpha,\beta\).
  Therefore the minimizer solves
  \[
    \min_{\alpha,\beta}\sum_g n_g(\bar Y_g-\alpha-\beta X_g)^2.
  \]
  \backbutton{grouped}
\end{frame}

\begin{frame}{Appendix A10: Controls and Precision}
  \hypertarget{app:precision}{}
  After residualizing treatment on controls, the treatment standard error has the form
  \[
    \mathrm{se}(\hat\tau)
    \approx
    \sqrt{\frac{\sigma^2}{\sum_i(\tilde T_i-\bar{\tilde T})^2}}.
  \]
  Adding a control can affect both terms:
  \begin{itemize}
    \item if it predicts \(Y\), residual variance \(\sigma^2\) may fall;
    \item if it predicts \(T\), residual treatment variance may fall;
    \item the standard error improves only if the first effect dominates the second.
  \end{itemize}
  \begin{block}{Takeaway}
    Outcome predictors are often useful precision controls. Pure treatment predictors can make the treatment coefficient noisier.
  \end{block}
  \backbutton{precision}
\end{frame}

\begin{frame}{Appendix A11: Mediators and Total Effects}
  \hypertarget{app:mediator}{}
  Suppose
  \[
    M_i=a+bT_i+v_i,
    \qquad
    Y_i=\alpha+\tau T_i+\theta M_i+\varepsilon_i.
  \]
  Substituting the mediator equation into the outcome:
  \[
    Y_i
    =
    (\alpha+\theta a)
    +
    (\tau+\theta b)T_i
    +
    \theta v_i+\varepsilon_i.
  \]
  The total effect is
  \[
    \tau+\theta b.
  \]
  If we control for \(M_i\), the coefficient on \(T_i\) is the direct effect \(\tau\), not the total effect.
  \vspace{0.3em}

  This may be useful for a direct-effect question, but it is wrong for a total-effect policy question.
  \backbutton{mediator}
\end{frame}

\begin{frame}{Appendix A12: Colliders and Induced Association}
  \hypertarget{app:collider}{}
  Let treatment be randomized:
  \[
    T_i\indep U_i,\qquad Y_i=\alpha+\tau T_i+\lambda U_i+\varepsilon_i.
  \]
  Now control for a common effect
  \[
    C_i=aT_i+bU_i+\eta_i.
  \]
  Before conditioning,
  \[
    \Cov(T_i,U_i)=0.
  \]
  But after conditioning on \(C_i\), high \(T_i\) must be partly offset by low \(U_i\) to keep \(C_i\) fixed.
  \vspace{0.3em}

  Thus, generally,
  \[
    \Cov(T_i,U_i\mid C_i)\neq 0.
  \]
  The regression coefficient on \(T_i\) then absorbs part of \(\lambda U_i\), even though \(T_i\) was randomized.
  \backbutton{collider}
\end{frame}

\end{document}
